\documentclass{../../presentation}

\title{Άλγεβρα - Ανισώσεις}
\subtitle{2ου βαθμού Ανισώσεις}
\author[Λόλας]{Κωνσταντίνος Λόλας }
\date{}

\begin{document}

\frame{\titlepage}

\begin{frame}{Ανισώσεις 2ου βαθμού}

  \begin{figure}
    \centering
    \includegraphics[width=0.8\textwidth]{images/butwhy.jpg}
  \end{figure}

\end{frame}

\begin{frame}{Εκτιμούμε τα παλιά}
  Αν μπορώ να λύσω ως προς $x$ έχω τελειώσει. Τα καινούρια...
  \begin{itemize}
    \item $x^2 -9> 0$
    \item $x^2-3x+2 \le 0$
    \item $5x^3-3x^2+2x-1 \ge 0$
  \end{itemize}
\end{frame}

\begin{frame}{Κάτι βαρετόοοοοοοοοοο}
  \begin{align*}
    αx^2+βx+γ & =α(x^2+\frac{β}{α}x+\frac{γ}{α})                                                 \\
              & =α(x^2+\frac{β}{α}x+\frac{β^2}{4α^2}-\frac{β^2}{4α^2}+\frac{γ}{α})               \\
              & = α\left( x^2+\frac{β}{α}x+\frac{β^2}{4α^2}-\frac{β^2}{4α^2}+\frac{γ}{α} \right) \\
              & = α\left(\left( x+\frac{β}{2α} \right)^2 - \frac{β^2-4αγ}{4α^2}\right)           \\
              & = α\left(\left( x+\frac{β}{2α} \right)^2 - \frac{Δ}{4α^2}\right)                 \\
  \end{align*}
\end{frame}

\begin{frame}{Περιπτώσεις}
  \begin{itemize}[<+->]
    \item $Δ>0$
          \begin{align*}
            α\left(\left( x+\frac{β}{2α} \right)^2 - \frac{Δ}{4α^2}\right) & =
            α\left( x+\frac{β}{2α} + \frac{\sqrt{Δ}}{2α} \right)\left( x+\frac{β}{2α} - \frac{\sqrt{Δ}}{2α} \right) \\
                                                                           & = α(x-ρ_1)(x-ρ_2)
          \end{align*}
    \item $Δ=0$
          \begin{align*}
            α\left(\left( x+\frac{β}{2α} \right)^2 - \frac{Δ}{4α^2}\right) & =
            α\left( x+\frac{β}{2α} \right) = α(x-ρ)^2
          \end{align*}
    \item $Δ<0$
          \begin{align*}
            α\left(\left( x+\frac{β}{2α} \right)^2 - \frac{Δ}{4α^2}\right) =
            α\left(+\right)
          \end{align*}
  \end{itemize}
\end{frame}

\begin{frame}{Παράδειγμα}

  \begin{enumerate}[<+->]
    \item $x^2-3x+2$
          $$x^2-3x+2=1(x-1)(x-2)$$
    \item $-2x^2+12x-18$
          $$-2(x^2-6x+9)=-2(x-3)^2$$
    \item $x^2+x+1$
          $$x^2+x+1=\left(x+\frac{1}{2}\right)^2+\frac{3}{4}$$
  \end{enumerate}

\end{frame}

\begin{frame}{Τι χρειαζόμαστε}
  Πρόσημο $x-ρ$
  \begin{center}
    \begin{tikzpicture}

      \tkzTabInit[lgt=2,espcl=1.5]
      {$x$ / 1 , $x-ρ$ / 1}
      {$-\infty$, $ρ$, $+\infty$}

      \tkzTabLine
      {,-,z ,+,}
    \end{tikzpicture}
  \end{center}

\end{frame}

\begin{frame}{Τι χρειαζόμαστε}
  Πρόσημο $(x-ρ_1)(x-ρ_2)?$
  \begin{center}
    \begin{tikzpicture}

      \tkzTabInit[lgt=4,espcl=1.5]
      {$x$ / 1 , $x-ρ_1$ / 1, $x-ρ_2$ / 1, $(x-ρ_1)(x-ρ_2)$ / 1}
      {$-\infty$, $ρ_1$, $ρ_2$, $+\infty$}

      \tkzTabLine
      {,-,z ,+,t,+,}
      \tkzTabLine
      {,-,t ,-,z,+,}
      \tkzTabLine
      {,+,z ,-,z,+,}

    \end{tikzpicture}
  \end{center}

\end{frame}

\begin{frame}{Αν $α<0$? }
  Πρόσημο $(x-ρ_1)(x-ρ_2)?$
  \begin{center}
    \begin{tikzpicture}

      \tkzTabInit[lgt=4,espcl=1.5]
      {$x$ / 1 , $α$ / 1, $x-ρ_1$ / 1, $x-ρ_2$ / 1, $α(x-ρ_1)(x-ρ_2)$ / 1}
      {$-\infty$, $ρ_1$, $ρ_2$, $+\infty$}

      \tkzTabLine
      {,-, t,-,t,-,}
      \tkzTabLine
      {,-,z ,+,t,+,}
      \tkzTabLine
      {,-,t ,-,z,+,}
      \tkzTabLine
      {,-,z ,+,z,-,}

    \end{tikzpicture}
  \end{center}

\end{frame}

\begin{frame}{Συγκεντρωτικά για $α>0$}
  \begin{center}
    \begin{itemize}
      \item $Δ > 0$
            \begin{center}
              \begin{tikzpicture}[scale=0.70]
                \tkzTabInit[lgt=4,espcl=1.5]
                {$x$/1, $ax^2+bx+c$/1}
                {$-\infty$, $ρ_1$, $ρ_2$, $+\infty$}
                \tkzTabLine
                {,+,z ,-,z,+,}
              \end{tikzpicture}
            \end{center}

      \item $Δ = 0$
            \begin{center}
              \begin{tikzpicture}[scale=0.70]
                \tkzTabInit[lgt=4,espcl=1.5]
                {$x$ / 1 , $ax^2+bx+c$ / 1}
                {$-\infty$, $ρ$, $+\infty$}

                \tkzTabLine
                {,+,z,+,}
              \end{tikzpicture}
            \end{center}

      \item $Δ < 0$
            \begin{center}
              \begin{tikzpicture}[scale=0.70]
                \tkzTabInit[lgt=4,espcl=1.5]
                {$x$ / 1 , $ax^2+bx+c$ / 1}
                {$-\infty$, $+\infty$}

                \tkzTabLine
                {,+,}
              \end{tikzpicture}
            \end{center}
    \end{itemize}

  \end{center}
\end{frame}

\begin{frame}{Συγκεντρωτικά για $α<0$}
  \begin{center}
    \begin{itemize}
      \item $Δ > 0$
            \begin{center}
              \begin{tikzpicture}[scale=0.70]
                \tkzTabInit[lgt=4,espcl=1.5]
                {$x$ / 1 , $ax^2+bx+c$ / 1}
                {$-\infty$, $ρ_1$, $ρ_2$, $+\infty$}
                \tkzTabLine
                {,-,z ,+,z,-,}
              \end{tikzpicture}
            \end{center}

      \item $Δ = 0$
            \begin{center}
              \begin{tikzpicture}[scale=0.70]
                \tkzTabInit[lgt=4,espcl=1.5]
                {$x$ / 1 , $ax^2+bx+c$ / 1}
                {$-\infty$, $ρ$, $+\infty$}

                \tkzTabLine
                {,-,z,-,}
              \end{tikzpicture}
            \end{center}

      \item $Δ < 0$
            \begin{center}
              \begin{tikzpicture}[scale=0.70]
                \tkzTabInit[lgt=4,espcl=1.5]
                {$x$ / 1 , $ax^2+bx+c$ / 1}
                {$-\infty$, $+\infty$}

                \tkzTabLine
                {,-,}
              \end{tikzpicture}
            \end{center}
    \end{itemize}

  \end{center}
\end{frame}

\section{ask}

\begin{askisi}
  Για τις διάφορες τιμές του $x\in\mathbb{R}$, να βρείτε το πρόσημο των τιμών των παρακάτω τριωνύμων.
  \begin{enumerate}[<+->]
    \item $x^2-3x+2$
    \item $-3x^2+x+2$
    \item $x^2-1$
    \item $-x^2+x$
  \end{enumerate}
\end{askisi}

\begin{askisi}
  Να κάνετε τον πίνακα προσήμων του τριωνύμου $x^2-2x+1$.
\end{askisi}

\begin{askisi}
  Να κάνετε τον πίνακα προσήμων των τριωνύμων
  \begin{enumerate}[<+->]
    \item $x^2-2x+2$
    \item $x-1-x^2$
    \item $-3x^2-1$
  \end{enumerate}
\end{askisi}

\begin{askisi}
  Να λύσετε την ανίσωση $x^2-x-2>0$.
\end{askisi}

\begin{askisi}
  Να λύσετε την ανίσωση $4x-2\ge x(3x-1)$.
\end{askisi}

\begin{askisi}
  Να λύσετε τις ανισώσεις
  \begin{enumerate}[<+->]
    \item $x^2\le 3x$
    \item $2x^2>1$
    \item $x^2-5x+6<0$
  \end{enumerate}
\end{askisi}

\begin{askisi}
  Να λύσετε τις ανισώσεις
  \begin{enumerate}[<+->]
    \item $4x^2>4x-1$
    \item $(x-1)^2>2x-4$
  \end{enumerate}
\end{askisi}

\begin{askisi}
  Να βρείτε τις τιμές του $x\in\mathbb{R}$ για τις οποίες συναληθεύουν οι ανισώσεις:
  $$x^2\le 9 \text{ και } x+2<x^2$$
\end{askisi}

\begin{askisi}
  Να βρείτε τις τιμές του $x\in\mathbb{R}$ για τις οποίες ισχύει:
  $$2x-1<x^2\le 3x$$
\end{askisi}

\begin{askisi}
  Δίνεται η εξίσωση $(λ-3)x^2-λx-1=0$, $λ\ne 3$. Να βρείτε τις τιμές του $λ$ για τις οποίες η εξίσωση έχει δύο ρίζες πραγματικές και άνισες.
\end{askisi}

\begin{askisi}
  Να βρείτε το πλήθος των ριζών της εξίσωσης $x^2-(2λ-1)x-2λ+1=0$ για κάθε τιμή του $λ$.
\end{askisi}

\begin{askisi}
  Να δείξετε ότι
  \begin{enumerate}[<+->]
    \item $λ^2-2λ+2>0$, για κάθε $λ\in\mathbb{R}$
    \item $x^2+xy+y^2>0$, για κάθε $x,y\in\mathbb{R}$
  \end{enumerate}
\end{askisi}

\begin{askisi}
  Να δείξετε ότι η εξίσωση $x^2+λx+λ-2=0$ έχει 2 πραγματικές ρίζες πραγματικές και άνισες για κάθε $λ\in\mathbb{R}$.
\end{askisi}

\begin{askisi}
  Να λύσετε την ανίσωση $x^2-3|x|+2<0$.
\end{askisi}


\begin{askisi}
  Να λύσετε την ανίσωση $|x^2-1|<x^2-x-2$.
\end{askisi}

\begin{askisi}
  Να βρείτε τις τιμές του $λ\in\mathbb{R}$ για τις οποίες ισχύει
  $$(λ-1)x^2-λx+λ<0 \text{, } λ\ne 1 \text{, για κάθε } x\in\mathbb{R}$$
\end{askisi}

\begin{askisi}
  Να βρείτε τις τιμές του $λ\in\mathbb{R}$ για τις οποίες ισχύει
  $$(λ-1)x^2+(λ-1)x+1<0 \text{, για κάθε } x\in\mathbb{R}$$
\end{askisi}

\begin{askisi}
  Για κάθε $x\in\mathbb{R}$ και $λ\in\mathbb{R}$, να αποδείξετε ότι $x^2-(λ+2)x+λ^2+2>0$.
\end{askisi}


\end{document}